\chapter{Kwantyfikatory i podstawowe metody dowodzenia}\label{ch:quantifiers}

\section{Dlaczego same spójniki nie wystarczają?}\label{sec:quantifier-motivation}
Zdania „$2$ jest parzyste” i „$7$ jest parzyste” mają
określone wartości logiczne. Matematyka posługuje się jednak
przede wszystkim wypowiedziami dotyczącymi całych rodzin
obiektów: „każda liczba parzysta ma pewną własność” albo
„istnieje liczba spełniająca równanie”. Samymi spójnikami
$\neg,\land,\lor,\Rightarrow,\Leftrightarrow$ nie umiemy
wygodnie zapisać, do \emph{ilu} obiektów odnosi się taki warunek.

Potrzebne są dwa nowe symbole: $\forall$ („dla każdego”)
oraz $\exists$ („istnieje”). W przeciwieństwie do zwykłych
spójników związane są one ze zmienną i ze wskazanym
\emph{zakresem jej wartości}.

\section{Dziedzina zmiennej i predykat}\label{sec:predicates}
\begin{definicja}[Dziedzina i predykat --- ujęcie wstępne]
\label{def:predicate}
Ustalamy pewną \emph{dziedzinę rozważań} $D$, czyli kolekcję
obiektów, o których chcemy mówić (może być także
pusta). \emph{Predykat} $P(x)$
to warunek o zmiennej $x$, który po wskazaniu obiektu
$x$ z $D$ przyjmuje wartość prawdy albo fałszu.
Samo $P(x)$ jest na ogół \emph{wypowiedzią otwartą},
a nie zdaniem o ustalonej wartości logicznej.
W tym rozdziale pojęcie dziedziny traktujemy intuicyjnie;
jego teoriomnogościową definicję podamy później.
\end{definicja}

\begin{uwaga}[Parametry i predykaty wieloargumentowe]
\label{rem:predicates-arity-xli}
Można również rozważać predykat $R(x,y)$ dwóch zmiennych
o dziedzinach odpowiednio $D$ i $E$: po wskazaniu
$a\in D$ i $b\in E$ otrzymujemy zdanie $R(a,b)$.
Przykładowo dla $D=E=\N$ predykat $R(x,y)$ może
oznaczać $x<y$. Napis $R(x,3)$ ma nadal wolną
zmienną $x$, natomiast $R(2,3)$ jest prawdziwym
zdaniem. Aby rozstrzygać wartość logiczną wyrażenia
z kwantyfikatorem i dodatkowymi zmiennymi wolnymi,
trzeba wcześniej ustalić ich wartości.
Ścisłą semantykę przypisań do zmiennych podamy w
rozdziale~\ref{ch:semantics-fol}.
\end{uwaga}

\begin{przyklad}[Ta sama wypowiedź, różna dziedzina]
Rozważmy $P(x)$: „$x^2=1$”. Jeśli dopuszczamy tylko
nieujemne liczby całkowite, jedynym rozwiązaniem jest $1$.
Jeśli dopuszczamy wszystkie liczby całkowite, rozwiązaniami
są $1$ oraz $-1$. Jeśli zaś dopuszczamy wszystkie liczby
rzeczywiste, nadal otrzymujemy te dwa rozwiązania.
Bez wskazania dziedziny wypowiedź „istnieje dokładnie jedno
rozwiązanie” nie jest dostatecznie określona.
\end{przyklad}

\begin{uwaga}[Podstawowa notacja liczbowa]
$\N=\{0,1,2,\ldots\}$ będzie oznaczać liczby naturalne,
$\Z$ liczby całkowite, a $\R$ liczby rzeczywiste.
Zapisy te wykorzystujemy na razie jako znane nazwy
kolekcji liczb. Dopiero w części teoriomnogościowej
uzasadnimy formalnie istnienie i konstrukcję takich zbiorów.
\end{uwaga}

\section{Kwantyfikator ogólny}\label{sec:forall}
\begin{definicja}[„Dla każdego”]\label{def:forall}
Zdanie $\forall x\in D\;P(x)$ jest prawdziwe wtedy i tylko
wtedy, gdy $P(a)$ jest prawdziwe dla \emph{każdego} obiektu
$a$ z dziedziny $D$. Jest fałszywe dokładnie wtedy, gdy
\emph{istnieje kontrprzykład}: pewien obiekt $a$ z $D$,
dla którego $P(a)$ jest fałszywe.
\end{definicja}

\begin{przyklad}[Wystarczy jeden kontrprzykład]
Zdanie $\forall n\in\N\;(n<10)$ jest fałszywe:
liczba $10$ jest naturalna, lecz nierówność $10<10$
jest fałszywa. Nie trzeba w tym celu analizować
pozostałych liczb naturalnych.
\end{przyklad}

\begin{przyklad}[Dlaczego trzeba ustalić dziedzinę?]
Wypowiedź „każde $x$ spełnia $x^2\ge0$” jest prawdziwa
dla liczb rzeczywistych. Przy przejściu do liczb zespolonych
nie staje się po prostu fałszywa: bez dodatkowego
określenia znaczenia symbolu $\ge$ nie ma ustalonej
wartości logicznej, ponieważ na liczbach zespolonych
nie przyjmujemy standardowo takiego porządku.
Nie wolno zmieniać dziedziny bez sprawdzenia
sensu użytych symboli.
\end{przyklad}

\section{Kwantyfikator egzystencjalny}\label{sec:exists}
\begin{definicja}[„Istnieje”]\label{def:exists}
Zdanie $\exists x\in D\;P(x)$ jest prawdziwe dokładnie
wtedy, gdy \emph{istnieje co najmniej jeden}
obiekt $a$ z dziedziny $D$, dla którego $P(a)$ jest
prawdziwe. Taki obiekt $a$ nazywamy \emph{świadkiem istnienia}.
Uwaga: prawdziwość zdania egzystencjalnego nie wymaga,
abyśmy potrafili obliczyć lub jawnie nazwać jego świadka.
W dowodach konstrukcyjnych go podajemy, ale w logice
klasycznej możliwe są także dowody niekonstrukcyjne.
Aby obalić zdanie egzystencjalne, trzeba wykazać,
że \emph{żaden} obiekt z $D$ nie spełnia warunku.
\end{definicja}

\begin{przyklad}[Wskazanie świadka]
Zdanie $\exists n\in\N\;(n^2=25)$ jest prawdziwe,
gdyż $n=5$ jest liczbą naturalną i $5^2=25$.
Zdanie $\exists n\in\N\;(n^2=-1)$ jest fałszywe,
ponieważ dla każdej liczby naturalnej $n$ zachodzi
$n^2\ge0$. W drugim przypadku pojedyncze
przykłady $n=0,1,2$ nie wystarczą do dowodu;
argument musi objąć wszystkie liczby naturalne.
\end{przyklad}

\begin{uwaga}[Dowód istnienia nie zawsze daje algorytm]
Nie należy utożsamiać warunku $\exists x\in D\;P(x)$
z istnieniem procedury, która dla danego $P$ wyszuka
świadka. Istnienie można czasem uzasadnić pośrednio,
wykazując sprzeczność założenia $\forall x\in D\;\neg P(x)$.
Jest to dowód klasyczny i nie musi wyznaczać konkretnego
$x$. W rozdziałach o obliczalności przekonamy się, dlaczego
takie rozróżnienie jest fundamentalne.
\end{uwaga}

\begin{przyklad}[Prawdziwość nie oznacza jednoznaczności]
Zdanie $\exists x\in\Z\;(x^2=1)$ jest prawdziwe,
ale nie mówi, że świadek jest jeden: spełniają je
$x=1$ i $x=-1$. Osobno wprowadzimy zapis
„istnieje dokładnie jeden”.
\end{przyklad}

\begin{uwaga}[Pusta dziedzina]
Jeżeli $D$ jest puste, to $\forall x\in D\;P(x)$ jest
prawdziwe: nie istnieje w $D$ obiekt obalający $P$.
Natomiast $\exists x\in D\;P(x)$ jest fałszywe:
nie można wskazać świadka w pustej dziedzinie.
Nie zakładamy przez to istnienia jakiegoś „pustego obiektu”.
To szczególne przypadki definicji obu kwantyfikatorów.
\end{uwaga}

\begin{uwaga}[Niepusta dziedzina struktur pierwszego rzędu]
\label{rem:empty-domain-semantics-xli}
W tym rozdziale dopuszczamy także pustą dziedzinę $D$
przy \emph{kwantyfikatorze ograniczonym}.
Nie należy z tego wyciągać wniosku, że dowolna
struktura klasycznej logiki pierwszego rzędu może
mieć puste uniwersum: w standardowo przyjętej
semantyce tej logiki uniwersum struktury jest niepuste.
Nie przeszkadza to interpretować $\forall x\in D$
i $\exists x\in D$ dla pustego \emph{podzbioru}
$D$ niepustego uniwersum, korzystając z rozwinięć
podanych w sekcji~\ref{sec:quantifier-scope}.
\end{uwaga}

\section{Zapisy skrócone i zakres kwantyfikatora}\label{sec:quantifier-scope}
\begin{definicja}[Zakres działania kwantyfikatora]
W zapisie $\forall x\in D\;(P(x)\Rightarrow Q(x))$
\emph{zakresem} kwantyfikatora jest cała formuła
w nawiasie. Wystąpienia zmiennej $x$ w tym zakresie są związane
przez dany kwantyfikator, o ile nie zostały związane
przez wewnętrzny kwantyfikator tej samej zmiennej.
Wystąpienie zmiennej, którego nie wiąże żaden
kwantyfikator, nazywamy \emph{wolnym}.
\end{definicja}

\begin{przyklad}[Zmienna wolna i związana]
Niech $y$ przebiega $\R$. W formule
$(\forall x\in\N\;x^2\ge0)\land(y>0)$
zmienna $x$ jest związana, natomiast $y$ pozostaje wolna:
wartość całego wyrażenia zależy jeszcze od wyboru $y$.
Jeśli dopiszemy $\exists y\in\R$ przed całą formułą,
uzyskamy zdanie bez zmiennych wolnych.
\end{przyklad}

\begin{uwaga}[Niebezpieczne nawiasy]
Formuły $\forall x\;(P(x)\Rightarrow Q(x))$
i $(\forall x\;P(x))\Rightarrow Q(x)$ nie wyrażają
tego samego: w drugiej formule wystąpienie $x$ w
$Q(x)$ jest wolne. Przy pierwszych obliczeniach
zawsze wpisujemy zakres kwantyfikatora w nawiasach.
\end{uwaga}

\begin{twierdzenie}[Zmiana nazwy zmiennej związanej]
\label{thm:alpha-renaming-xli}
Niech $D$ będzie ustaloną dziedziną, a $P(x,z)$
predykatem, w którym $z$ oznacza ewentualne dodatkowe
parametry. Niech $y$ będzie nową nazwą zmiennej,
która nie występuje w $P(x,z)$. Wtedy, po ustaleniu
wartości parametrów $z$, zdania
$\forall x\in D\;P(x,z)$ i
$\forall y\in D\;P(y,z)$ są równoważne; podobnie
$\exists x\in D\;P(x,z)$ i
$\exists y\in D\;P(y,z)$ są równoważne.
W drugich formułach $P(y,z)$ oznacza wynik zastąpienia
\emph{wolnych} wystąpień $x$ przez $y$.
\end{twierdzenie}
\begin{proof}
Ustalmy wartości wszystkich parametrów $z$.
Dla każdego obiektu $a\in D$ zdania otrzymane
przez podstawienie $x=a$ w $P(x,z)$ i przez
podstawienie $y=a$ w $P(y,z)$ mają tę samą treść:
zmieniliśmy nazwę miejsca podstawienia, a nie
sam warunek na $a$. Zdanie ogólne po obu stronach
wymaga zatem prawdziwości tej samej rodziny zdań
$P(a,z)$ dla wszystkich $a\in D$.
Zdanie egzystencjalne po obu stronach wymaga
istnienia świadka w tej samej rodzinie.
W obu przypadkach warunki prawdziwości są identyczne.
\end{proof}

\begin{przyklad}[Niebezpieczeństwo przechwycenia zmiennej]
\label{ex:variable-capture-xli}
Przemianowanie $x$ na nazwę występującą pod innym
kwantyfikatorem wymaga sprawdzenia, czy nie dojdzie
do przechwycenia zmiennej. Na przykład
$\forall x\in\N\;\exists y\in\N\;(x<y)$ jest prawdziwe:
dla każdego $x$ wybieramy $y=x+1$.
Błędne przemianowanie zewnętrznej zmiennej $x$ na $y$
(łącznie z jej wystąpieniem w nierówności) prowadzi do
$\forall y\in\N\;\exists y\in\N\;(y<y)$;
wewnętrzny kwantyfikator wiąże teraz oba wystąpienia
$y$ w nierówności, która jest zawsze fałszywa.
Poprawna zmiana obu nazw z zachowaniem ich odrębności
daje np.
$\forall u\in\N\;\exists v\in\N\;(u<v)$.
\end{przyklad}

\begin{przyklad}[Wszystkie elementy spełniające warunek]
Wypowiedź „każda liczba naturalna podzielna przez $4$
jest parzysta” zapisujemy
\[
\forall n\in\N\;\bigl(4\mid n\Rightarrow2\mid n\bigr),
\]
gdzie $a\mid b$ czytamy „$a$ dzieli $b$”.
Zapis $4\mid n$ znaczy, że $n=4k$ dla pewnej
liczby całkowitej $k$. Gdy $4$ nie dzieli $n$,
nie naruszamy prawdziwości implikacji.
\end{przyklad}

\begin{uwaga}[Skrót formalny]
Załóżmy, że $D$ jest podzbiorem ustalonego niepustego
uniwersum $U$, a warunek $P(x)$ ma określoną wartość
logiczną dla każdego $x\in U$. Wtedy zapis
$\forall x\in D\;P(x)$ jest skrótem formuły
$\forall x\;(x\in D\Rightarrow P(x))$;
podobnie $\exists x\in D\;P(x)$ oznacza
$\exists x\;(x\in D\land P(x))$.
Bez takiego wspólnego uniwersum i interpretacji $P$
poza $D$ nie musimy mieć prawa używać tych rozwinięć:
poprzednie definicje kwantyfikatorów ograniczonych
wymagały wartości $P(x)$ jedynie dla $x\in D$.
Symbol $\in$ będzie zdefiniowany dokładnie w następnym
rozdziale; na razie oznacza po prostu „jest elementem”.
\end{uwaga}

\begin{uwaga}[Kiedy rozwinięcia skrótów są prawidłowe?]
W zapisie $\forall x\,(x\in D\Rightarrow P(x))$ zmienna
$x$ przebiega całe ustalone uniwersum, a $x\in D$ jest
dodatkowym warunkiem. Jeśli $D$ jest puste, poprzednik
implikacji jest zawsze fałszywy, zatem kwantyfikator
ogólny jest prawdziwy. W zapisie
$\exists x\,(x\in D\land P(x))$ koniunkcja nie może
być prawdziwa dla żadnego $x$.
Właśnie dlatego w ograniczeniu kwantyfikatora ogólnego
występuje \emph{implikacja}, a w ograniczeniu kwantyfikatora
egzystencjalnego --- \emph{koniunkcja}; ich zamiana
zmieniłaby sens sformułowania.
\end{uwaga}

\begin{przyklad}[Jak odróżnić „nie wszyscy” od „nikt”?]
Weźmy wypowiedź „wszyscy uczniowie zdali egzamin”.
Jej negacją jest „istnieje uczeń, który nie zdał”, a nie
„żaden uczeń nie zdał”. Pierwsze zdanie kwestionuje choć jeden
przypadek; drugie kwestionuje wszystkie. Analogicznie negacja
„istnieje uczeń, który zdał” brzmi „żaden uczeń nie zdał”.
To przygotowuje nas do ścisłych wzorów z kwantyfikatorami.
\end{przyklad}

\section{Negowanie kwantyfikatorów}\label{sec:quantifier-negation}
\begin{twierdzenie}[Prawa de Morgana dla kwantyfikatorów]
\label{thm:negation-quantifiers}
Niech $D$ będzie dowolną dziedziną (także pustą), a
$P(x)$ predykatem o zmiennej $x$ z dziedziny $D$.
W klasycznej logice zachodzą równoważności
\[
\neg\bigl(\forall x\in D\;P(x)\bigr)
\ \Leftrightarrow\ \exists x\in D\;\neg P(x),
\]
\[
\neg\bigl(\exists x\in D\;P(x)\bigr)
\ \Leftrightarrow\ \forall x\in D\;\neg P(x).
\]
\end{twierdzenie}
\begin{proof}
Pierwsza równoważność: zdanie $\forall x\in D\;P(x)$
jest fałszywe dokładnie wtedy, gdy istnieje w $D$
co najmniej jeden kontrprzykład dla $P$, czyli
obiekt $a$ taki, że $P(a)$ jest fałszywe.
Na mocy definicji negacji jest to dokładnie warunek,
że $\neg P(a)$ jest prawdziwe, a więc że
$\exists x\in D\;\neg P(x)$ zachodzi.

Druga równoważność: $\exists x\in D\;P(x)$
jest fałszywe dokładnie wtedy, gdy nie ma w $D$
żadnego świadka dla $P$. To znaczy, że dla
każdego obiektu $a$ z $D$ formuła $P(a)$
jest fałszywa, czyli $\neg P(a)$ jest prawdziwa.
Definicja kwantyfikatora ogólnego daje
$\forall x\in D\;\neg P(x)$.
W obu dowodach przechodziliśmy przez równoważne
warunki, więc uzyskaliśmy obie strony implikacji.
\end{proof}

\begin{przyklad}[Negowanie zdania z codzienną treścią]
„Każdy student zdał egzamin” negujemy jako
„Istnieje student, który nie zdał egzaminu”.
Nie jako „Żaden student nie zdał egzaminu”:
ta druga wypowiedź jest znacznie silniejsza.
Podobnie zaprzeczenie „Istnieje rozwiązanie równania”
brzmi „Żaden dopuszczalny obiekt nie jest rozwiązaniem”,
a nie „Istnieje obiekt, który nie jest rozwiązaniem”.
\end{przyklad}

\begin{przyklad}[Negacja implikacji pod kwantyfikatorem]
Rozważmy zdanie
$\forall n\in\N\;(n>2\Rightarrow n^2>4)$.
Jego zaprzeczenie ma postać
$\exists n\in\N\;(n>2\land n^2\le4)$.
Dlaczego koniunkcja? Ponieważ
$\neg(p\Rightarrow q)\equiv p\land\neg q$
(twierdzenie~\ref{thm:negated-implication-xli}):
aby obalić implikację, musimy mieć zarazem
prawdziwy poprzednik i fałszywy następnik.
\end{przyklad}

\section{Kolejność kwantyfikatorów ma znaczenie}
\label{sec:quantifier-order}
\begin{przyklad}[„Każdemu można dobrać” a „istnieje jeden dla wszystkich”]
Rozpatrzmy na liczbach naturalnych warunek $x<y$.
Zdanie
\[
\forall x\in\N\;\exists y\in\N\;(x<y)
\]
jest prawdziwe: dla dowolnie wskazanego $x$
możemy wybrać $y=x+1$. Wybrany $y$ \emph{zależy od $x$}.
Natomiast zdanie
\[
\exists y\in\N\;\forall x\in\N\;(x<y)
\]
jest fałszywe. Gdyby istniało takie $y$,
można byłoby wstawić $x=y$, otrzymując $y<y$,
co jest fałszem. W szczególności nie istnieje
jedna liczba naturalna większa od wszystkich
liczb naturalnych.
\end{przyklad}

\begin{twierdzenie}[Podwójne kwantyfikatory tego samego typu]
\label{thm:same-quantifier-order}
Niech $D,E$ będą dziedzinami, a $P(x,y)$ predykatem
dla $x\in D$ i $y\in E$. Wtedy
\[
\forall x\in D\;\forall y\in E\;P(x,y)
\ \Leftrightarrow\
\forall y\in E\;\forall x\in D\;P(x,y),
\]
\[
\exists x\in D\;\exists y\in E\;P(x,y)
\ \Leftrightarrow\
\exists y\in E\;\exists x\in D\;P(x,y).
\]
\end{twierdzenie}
\begin{proof}
Pierwsze zdanie po lewej jest prawdziwe dokładnie
wtedy, gdy dla \emph{każdej} pary dopuszczalnych
obiektów $a\in D$, $b\in E$ warunek $P(a,b)$
jest prawdziwy. Kolejność, w jakiej wymienimy te
dwa obiekty, nie zmienia wartości $P(a,b)$.
To samo kryterium charakteryzuje zdanie po prawej.

Druga para zdań ma po obu stronach dokładnie
ten sam warunek: można wskazać \emph{co najmniej
jedną} parę $a\in D$, $b\in E$, dla której
$P(a,b)$ jest prawdziwe. Wskazanie najpierw $a$,
a potem $b$ albo odwrotnie nie zmienia
istnienia takiej pary.
\end{proof}

\begin{twierdzenie}[Jedna poprawna zamiana kwantyfikatorów różnych typów]
\label{thm:swap-mixed}
Niech $D,E$ będą dziedzinami, a $P(x,y)$ predykatem
określonym dla $x\in D$ i $y\in E$. Wtedy
\[
\Bigl(\exists y\in E\;\forall x\in D\;P(x,y)\Bigr)
\Rightarrow
\Bigl(\forall x\in D\;\exists y\in E\;P(x,y)\Bigr).
\]
Implikacja odwrotna nie zachodzi ogólnie.
\end{twierdzenie}
\begin{proof}
Załóżmy, że po lewej istnieje obiekt $b\in E$,
dla którego $P(a,b)$ zachodzi dla każdego $a\in D$.
Weźmy dowolny $a\in D$. Ten sam, już wybrany
obiekt $b\in E$ spełnia $P(a,b)$, a więc
istnieje odpowiedni $y\in E$. Ponieważ wybór
$a\in D$ był dowolny, otrzymujemy zdanie po prawej.
Jeżeli dziedzina $D$ jest pusta, wniosek po prawej jest
prawdziwy niezależnie od poprzednika, więc
implikacja również pozostaje prawdziwa.

Aby pokazać, że odwrotność jest fałszywa,
weźmy $D=E=\N$ i $P(x,y)$ określony przez $x<y$.
Jak wykazaliśmy w poprzednim przykładzie,
$\forall x\in\N\;\exists y\in\N\;(x<y)$
jest prawdziwe, a
$\exists y\in\N\;\forall x\in\N\;(x<y)$
jest fałszywe. Jest to kontrprzykład do odwrotności.
\end{proof}

\begin{uwaga}[Puste dziedziny w zamianie kwantyfikatorów]
\label{rem:mixed-quantifiers-empty-domains-lxiii}
W twierdzeniu~\ref{thm:swap-mixed} poprawność implikacji
nie wymaga niepustości $D$ ani $E$, o ile dopuszczamy
semantykę kwantyfikatorów po pustych zbiorach.
Jeżeli $D=\varnothing$, to zdanie
$\forall x\in D\;\exists y\in E\;P(x,y)$ jest prawdziwe,
niezależnie od $E$; zdanie
$\exists y\in E\;\forall x\in D\;P(x,y)$
jest natomiast prawdziwe dokładnie wtedy, gdy $E\ne\varnothing$.
Jeżeli $D\ne\varnothing$ i $E=\varnothing$, obie strony są
fałszywe, ponieważ nie ma świadka $y$.
Wszystkie te przypadki zachowują prawdziwość implikacji,
ale pokazują, że nie można bez dodatkowych założeń
zamienić jej w równoważność.
\end{uwaga}
\begin{przyklad}[Kontrprzykład o pustym zbiorze świadków]
\label{ex:mixed-quantifiers-empty-lxiii}
Weźmy $D=\varnothing$ i $E=\varnothing$.
Wtedy $\forall x\in D\;\exists y\in E\;P(x,y)$ jest
prawdziwe na mocy pustego zakresu kwantyfikatora ogólnego,
a $\exists y\in E\;\forall x\in D\;P(x,y)$ jest fałszywe,
ponieważ kwantyfikator egzystencjalny nie ma świadka.
Jest to kontrprzykład do implikacji odwrotnej niezależny
od arytmetyki liczb naturalnych. W standardowej semantyce
pierwszego rzędu nośnik całej struktury jest niepusty;
tutaj $D$ i $E$ są zakresami kwantyfikatorów ograniczonych,
które mogą być puste.
\end{przyklad}

\begin{przyklad}[Negowanie dwóch kwantyfikatorów]
Zaprzeczeniem
$\forall x\in\R\;\exists y\in\R\;(x+y=0)$
jest
$\exists x\in\R\;\forall y\in\R\;(x+y\ne0)$.
Przenosimy negację kolejno przez oba kwantyfikatory,
każdorazowo zmieniając ich typ, aż trafi do
wewnętrznego warunku. Oryginalne zdanie jest
prawdziwe: wystarczy dla danego $x$ wybrać $y=-x$.
Jego negacja jest zatem fałszywa.
\end{przyklad}

\section{Istnienie dokładnie jednego obiektu}\label{sec:unique-existence}
\begin{definicja}[Kwantyfikator jednoznacznego istnienia]
\label{def:unique}
Zapis $\exists!x\in D\;P(x)$ czytamy:
„istnieje \emph{dokładnie jeden} obiekt $x$ z $D$,
dla którego $P(x)$ jest prawdziwe”. Oznacza on
łącznie dwa warunki:
\begin{enumerate}[label=(\roman*)]
\item \emph{istnienie}: pewien $a\in D$ spełnia $P(a)$;
\item \emph{jednoznaczność}: jeśli $a,b\in D$ spełniają
      $P$, to $a=b$.
\end{enumerate}
Nie jest to nowy podstawowy kwantyfikator: jest to
wygodny skrót łączący $\exists$ i $\forall$.
\end{definicja}

\begin{twierdzenie}[Rozwinięcie symbolu $\exists!$]
\label{thm:unique-expansion-xli}
Dla dowolnej dziedziny $D$ (także pustej) i predykatu
$P(x)$ następujące zdania są równoważne:
\begin{align*}
&\exists!x\in D\;P(x),\\
&\exists x\in D\;\Bigl(P(x)\land
  \forall y\in D\,(P(y)\Rightarrow y=x)\Bigr).
\end{align*}
\end{twierdzenie}
\begin{proof}
Załóżmy, że pierwsze zdanie jest prawdziwe.
Z definicji istnieje świadek $a\in D$ spełniający $P(a)$.
Jeżeli dowolny $b\in D$ spełnia $P(b)$, jednoznaczność
wymusza $b=a$. Zatem dla $x=a$ prawdziwa jest
koniunkcja w zakresie kwantyfikatora egzystencjalnego
w drugim zdaniu, więc to zdanie zachodzi.

Odwrotnie, niech drugie zdanie będzie prawdziwe
i niech $a\in D$ będzie jego świadkiem.
Pierwszy człon koniunkcji daje $P(a)$, czyli istnienie.
Aby wykazać jednoznaczność, weźmy dowolne $b\in D$
spełniające $P(b)$. Drugi człon koniunkcji zastosowany
do $y=b$ daje $P(b)\Rightarrow b=a$, zatem $b=a$.
To jest dokładnie definicja jednoznacznego istnienia.
Jeżeli $D$ jest puste, oba zdania są fałszywe,
ponieważ w obu musi wystąpić świadek $x\in D$.
\end{proof}

\begin{przyklad}[Jednoznaczne rozwiązanie równania]
Zdanie $\exists!x\in\R\;(2x+3=7)$ jest prawdziwe.
Istnienie: $x=2$ spełnia równość, bo $2\cdot2+3=7$.
Jednoznaczność: jeśli $2x+3=7$, odejmujemy $3$ od
obu stron i dostajemy $2x=4$; dzieląc przez $2$,
uzyskujemy $x=2$. Każde rozwiązanie musi więc
być tym samym obiektem.
\end{przyklad}

\begin{przyklad}[Istnienie bez jednoznaczności]
$\exists x\in\Z\;(x^2=4)$ jest prawdziwe,
lecz $\exists!x\in\Z\;(x^2=4)$ jest fałszywe,
ponieważ $2^2=4$ i $(-2)^2=4$, a $2\ne-2$.
\end{przyklad}

\section{Jak budować matematyczny dowód?}\label{sec:proof-methods}
\paragraph{Co oznacza udowodnić twierdzenie?}
Twierdzenie ma \emph{tezę} i może mieć \emph{założenia}.
Dowód to skończony ciąg uzasadnionych przejść
od przyjętych założeń i wcześniej ustalonych faktów
do tezy. W dalszych partiach książki zdefiniujemy
formalny rachunek dowodowy, ale już teraz możemy
nauczyć się jego najważniejszych poprawnych schematów.

\subsection{Dowód bezpośredni}\label{sec:direct-proof}
Aby wykazać $p\Rightarrow q$, zakładamy $p$
i na podstawie definicji oraz znanych twierdzeń
wyprowadzamy $q$. Nie wystarczy sprawdzić
kilku przykładów spełniających $p$.

\begin{przyklad}[Suma dwóch liczb parzystych]
\label{ex:even-sum}
\emph{Teza.} Dla dowolnych liczb całkowitych $a,b$,
jeżeli obie są parzyste, to ich suma jest parzysta.

\emph{Dowód.} Z definicji parzystości $a=2k$
i $b=2l$ dla pewnych liczb całkowitych $k,l$.
Dodając równości, otrzymujemy
$a+b=2k+2l=2(k+l)$. Suma $k+l$ jest liczbą
całkowitą, więc $a+b$ ma postać dwukrotności
liczby całkowitej. Z definicji jest zatem parzysta.
\end{przyklad}

\begin{uwaga}[Parzyste i nieparzyste liczby całkowite]
Korzystamy z elementarnego faktu arytmetycznego, że każda
liczba całkowita $n$ ma dokładnie jedną z postaci
$n=2k$ oraz $n=2k+1$, gdzie $k\in\Z$.
Jest to szczególny przypadek dzielenia z resztą przez $2$:
reszta jest równa $0$ albo $1$. Pierwsza postać definiuje
parzystość, druga nieparzystość. Nie zachodzą jednocześnie,
ponieważ równość $2k=2l+1$ prowadziłaby do
$2(k-l)=1$, co jest niemożliwe dla liczb całkowitych.
Formalny dowód istnienia reszty z dzielenia podamy przy
budowaniu arytmetyki liczb naturalnych; nie wyprowadzamy
go z samych praw logiki.
\end{uwaga}

\subsection{Dowód przez przeciwstawienie}\label{sec:contrapositive-proof}
Z rozdziału~\ref{ch:propositional} wiemy, że
$p\Rightarrow q$ jest równoważne
$\neg q\Rightarrow\neg p$. Czasem znacznie
łatwiej udowodnić drugi z tych warunków.

\begin{przyklad}[Jeśli kwadrat jest parzysty, liczba jest parzysta]
\label{ex:even-square}
\emph{Teza.} Dla każdego $n\in\Z$, jeśli $n^2$
jest parzyste, to $n$ jest parzyste.

\emph{Dowód.} Udowodnimy implikację przeciwstawną.
Załóżmy, że $n$ jest \emph{nieparzyste}. Wtedy
$n=2k+1$ dla pewnego $k\in\Z$. Podnosimy
do kwadratu:
\[
n^2=(2k+1)^2=4k^2+4k+1=2(2k^2+2k)+1.
\]
Ponieważ $2k^2+2k$ jest liczbą całkowitą,
$n^2$ ma postać $2m+1$, a więc jest nieparzyste.
Pokazaliśmy, że z nieparzystości $n$ wynika
nieparzystość $n^2$. Na mocy równoważności
implikacji z przeciwstawną wynika stąd teza.
\end{przyklad}

\subsection{Dowód nie wprost}\label{sec:contradiction-proof}
Aby udowodnić zdanie $q$, możemy założyć jego
negację $\neg q$ i wyprowadzić sprzeczność:
jednoczesną prawdziwość pewnego zdania $r$
i jego negacji. W klasycznej logice oznacza to,
że $\neg q$ nie może być prawdziwe, więc $q$
musi być prawdziwe.

\begin{przyklad}[Niewymierność pierwiastka z dwóch]
\label{ex:sqrt-two}
Przypomnijmy znaczenie słów: liczba rzeczywista $t$
jest \emph{wymierna}, jeżeli $t=a/b$ dla pewnych
liczb całkowitych $a,b$ z $b\ne0$; w przeciwnym
razie jest \emph{niewymierna}. Przez $\sqrt2$
rozumiemy dodatnią liczbę rzeczywistą, której
kwadrat jest równy $2$. Wykorzystamy elementarne
twierdzenie arytmetyki, że ułamek można skrócić
do postaci, w której licznik i mianownik nie mają
wspólnego dodatniego dzielnika większego od $1$.
Formalne uzasadnienie tego faktu zostanie podane
w części poświęconej konstrukcji liczb naturalnych.

\emph{Teza.} Liczba $\sqrt2$ jest niewymierna.

\emph{Dowód.} Przypuśćmy przeciwnie, że $\sqrt2$
jest wymierna. Istnieją więc liczby całkowite
$a,b$ z $b\ne0$ i $\sqrt2=a/b$. Po skróceniu
ułamka możemy przyjąć, że $a$ i $b$ nie mają
wspólnego dzielnika większego od $1$; to znaczy,
że ułamek jest nieskracalny. Podnosząc równość
do kwadratu i mnożąc przez $b^2$, otrzymujemy
$a^2=2b^2$. Zatem $a^2$ jest parzyste, więc
$a$ jest parzyste na mocy Przykładu~\ref{ex:even-square}.
Możemy więc zapisać $a=2k$ dla pewnego $k\in\Z$.
Podstawiając do poprzedniej równości, mamy
$4k^2=2b^2$, a po podzieleniu przez $2$:
$b^2=2k^2$. Także $b^2$ jest parzyste, więc
$b$ jest parzyste. Wynika stąd, że zarówno $a$,
jak i $b$ są podzielne przez $2$, co przeczy
nieskracalności ułamka $a/b$. Otrzymana
sprzeczność obala przypuszczenie o wymierności
$\sqrt2$; dowodzi to tezy.
\end{przyklad}

\begin{uwaga}[Gdzie użyliśmy logiki klasycznej?]
Przejście „założenie $\neg q$ prowadzi do sprzeczności,
zatem $q$” korzysta z zasady usuwania podwójnej
negacji. W logice intuicjonistycznej trzeba staranniej
rozróżniać dowód $\neg\neg q$ od dowodu $q$.
W tej książce na razie pracujemy w logice klasycznej.
\end{uwaga}

\subsection{Dowód przez przypadki}\label{sec:proof-by-cases}
Jeżeli wiadomo, że zachodzi $p\lor q$, można
osobno wykazać tezę $r$ przy założeniu $p$,
a następnie przy założeniu $q$. Oznacza to,
że bez względu na to, która z dopuszczalnych
możliwości zachodzi, otrzymujemy $r$.

\begin{przyklad}[Iloczyn dwóch kolejnych liczb jest parzysty]
\emph{Teza.} Dla każdej liczby całkowitej $n$
iloczyn $n(n+1)$ jest parzysty.

\emph{Dowód.} Każda liczba całkowita jest parzysta
albo nieparzysta. W pierwszym przypadku $n=2k$
dla pewnego $k\in\Z$, więc
$n(n+1)=2[k(n+1)]$ jest parzyste. W drugim
przypadku $n=2k+1$, zatem $n+1=2(k+1)$
i $n(n+1)=2[n(k+1)]$ jest parzyste.
Oba przypadki wyczerpują możliwości, więc
teza zachodzi dla każdego $n\in\Z$.
\end{przyklad}

\subsection{Jak dowodzić równoważności i jednoznaczności?}
\label{sec:proof-biconditional}
Aby uzasadnić $p\Leftrightarrow q$, udowadniamy
osobno $p\Rightarrow q$ i $q\Rightarrow p$.
Aby uzasadnić $\exists!x\in D\;P(x)$ dla ustalonej
dziedziny $D$, wykazujemy istnienie obiektu z $D$
spełniającego $P$ (na przykład wskazując świadka),
a następnie porównujemy dowolne dwa obiekty z $D$
spełniające $P$ i wykazujemy, że muszą być równe.
Dowód samego istnienia nie uzasadnia jednoznaczności.

\begin{przyklad}[Równoważność dla parzystości]
\emph{Teza.} Dla każdej liczby całkowitej $n$
liczba $n$ jest parzysta wtedy i tylko wtedy,
gdy liczba $n+2$ jest parzysta.

\emph{Dowód.} Jeśli liczba $n$ jest parzysta, to $n=2k$
dla pewnego $k\in\Z$, a więc $n+2=2(k+1)$,
czyli liczba $n+2$ jest parzysta. Odwrotnie,
jeżeli liczba $n+2$ jest parzysta, to $n+2=2l$
dla pewnego $l\in\Z$. Po odjęciu $2$
uzyskujemy $n=2(l-1)$, więc również liczba $n$
jest parzysta. Wykazaliśmy oba kierunki.
\end{przyklad}

\section{Indukcja matematyczna --- pierwsze zastosowanie}
\label{sec:induction}
Liczb naturalnych jest nieskończenie wiele.
Nie da się zakończyć skończonego dowodu przez oddzielne
sprawdzanie wszystkich warunków $P(0),P(1),P(2),\ldots$.
Indukcja jest metodą,
która łączy sprawdzenie jednego przypadku z
regułą przejścia do następnego. Jej formalne
podstawy, wraz z konstrukcją liczb naturalnych,
rozwiną późniejsze rozdziały.

\begin{twierdzenie}[Zasada indukcji --- postać elementarna]
\label{thm:induction}
Niech $P(n)$ będzie predykatem dla $n\in\N$.
Jeśli zachodzą dwa warunki:
\begin{enumerate}[label=(\roman*)]
\item $P(0)$ jest prawdziwe (\emph{baza indukcji});
\item dla każdego $n\in\N$ z prawdziwości $P(n)$
      wynika prawdziwość $P(n+1)$ (\emph{krok indukcyjny}),
\end{enumerate}
to $P(n)$ jest prawdziwe dla każdego $n\in\N$.
\end{twierdzenie}

\begin{uwaga}[Status zasady indukcji]
Na tym etapie traktujemy zasadę indukcji jako
podstawową własność liczb naturalnych, a nie jako
wniosek z samych tablic prawdy. Później pokażemy,
jak zasada ta wiąże się z budową $\N$ i aksjomatami
arytmetyki. Nie wolno uznać jej za „oczywisty”
wniosek z faktu, że sprawdziliśmy kilka liczb.
\end{uwaga}

\begin{przyklad}[Wzór na sumę pierwszych liczb naturalnych]
\label{ex:gauss-induction}
\emph{Teza.} Dla każdego $n\in\N$ zachodzi
\[
0+1+\cdots+n=\frac{n(n+1)}2.
\]
\emph{Dowód.} Dla $n=0$ lewa strona wynosi $0$,
a prawa $0\cdot1/2=0$; baza zachodzi.

Ustalmy teraz dowolne $n\in\N$ i załóżmy,
że wzór jest prawdziwy dla tej liczby:
$0+1+\cdots+n=n(n+1)/2$. Chcemy wykazać
wzór dla $n+1$. Dodajemy $n+1$ do obu stron
równości przyjętej w założeniu indukcyjnym:
\begin{align*}
0+1+\cdots+n+(n+1)
&=\frac{n(n+1)}2+(n+1)\\
&=\frac{n(n+1)+2(n+1)}2\\
&=\frac{(n+1)(n+2)}2.
\end{align*}
Ostatni wyraz jest dokładnie prawą stroną
wzoru po zastąpieniu $n$ przez $n+1$.
Udowodniliśmy krok indukcyjny, więc
na mocy zasady indukcji wzór obowiązuje
dla wszystkich liczb naturalnych.
\end{przyklad}

\begin{przyklad}[Suma kolejnych liczb nieparzystych]
\label{ex:odd-sum-induction-xli}
Dla każdego $n\in\N$ zachodzi
$1+3+\cdots+(2n+1)=(n+1)^2$; po lewej jest
$n+1$ kolejnych liczb nieparzystych.
Dla $n=0$ lewa strona wynosi $1$, a prawa $(0+1)^2=1$.
Przyjmijmy teraz dla ustalonego $n\in\N$, że
$1+3+\cdots+(2n+1)=(n+1)^2$. Aby otrzymać wzór
dla $n+1$, trzeba dodać następny składnik, czyli
$2(n+1)+1=2n+3$. Z założenia indukcyjnego:
\[
1+3+\cdots+(2n+1)+(2n+3)
 =(n+1)^2+2n+3=(n+2)^2.
\]
Jest to dokładnie prawa strona wzoru dla $n+1$.
Warunki zasady indukcji są spełnione, więc
równość zachodzi dla każdego $n\in\N$.
\end{przyklad}

\begin{przyklad}[Dlaczego dowód dla jednego $n$ nie stanowi kroku indukcyjnego]
Załóżmy, że chcemy udowodnić $P(n)$ dla wszystkich
$n\in\N$. Sprawdzenie $P(0)$ i $P(1)$ nie wystarcza;
potrzebujemy zdania $\forall n\in\N\,(P(n)\Rightarrow P(n+1))$.
W kroku wybieramy \emph{dowolne} $n$, zakładamy dla niego
$P(n)$, a następnie wyprowadzamy $P(n+1)$.
Nie można zakładać $P(n+1)$ --- to właśnie teza kroku.
\end{przyklad}

\begin{przyklad}[Fałszywy krok indukcyjny ukryty w przykładach]
\label{ex:induction-counterexample-xli}
Rozważmy $P(n)$: „$n<3$” dla $n\in\N$.
Zdania $P(0)$, $P(1)$ i $P(2)$ są prawdziwe.
Co więcej, implikacje $P(0)\Rightarrow P(1)$
i $P(1)\Rightarrow P(2)$ są prawdziwe.
Jednak $P(3)$ jest fałszywe. Dokładnie dla $n=2$
implikacja $P(2)\Rightarrow P(3)$ jest fałszywa:
jej poprzednik ma wartość $1$, a następnik $0$.
Zatem z poprawnej bazy i kilku sprawdzonych
przejść nie wynika jeszcze warunek
$\forall n\in\N\,(P(n)\Rightarrow P(n+1))$.
\end{przyklad}

\section{Typowe pomyłki przy dowodzeniu}\label{sec:proof-errors}
\begin{przyklad}[Wstawienie liczby nie dowodzi zdania ogólnego]
Dla $n=2$ zachodzi $n^2\ge n$, ale nie dowodzi
to jeszcze twierdzenia
$\forall n\in\N\;(n^2\ge n)$.
Aby je wykazać, trzeba wyprowadzić nierówność
bez ograniczania się do jednego $n$.
Na przykład dla $n=0$ obie strony są równe,
a dla $n\ge1$ mamy $n^2-n=n(n-1)\ge0$.
Ta argumentacja rzeczywiście obejmuje każdy
możliwy naturalny $n$.
\end{przyklad}

\begin{przyklad}[Wybór świadka zależnego od zmiennej]
W prawdziwym zdaniu
$\forall x\in\R\;\exists y\in\R\;(x+y=0)$
wybrany świadek $y=-x$ zależy od rozważanego $x$.
Nie uprawnia to do zapisania
$\exists y\in\R\;\forall x\in\R\;(x+y=0)$.
Takie zdanie wymagałoby jednego ustalonego $y$,
działającego jednocześnie dla wszystkich $x$.
Dla $x=0$ wymuszałoby $y=0$, a dla $x=1$
wymuszałoby $y=-1$, co jest niemożliwe.
\end{przyklad}

\begin{przyklad}[Założenie tezy nie jest dowodem]
Chcąc wykazać $p\Rightarrow q$, wolno przyjąć
$p$, ale nie wolno bez uzasadnienia przyjąć $q$.
Z kolei w dowodzie nie wprost można przyjąć
$\neg q$, ale wtedy trzeba rzeczywiście
wyprowadzić sprzeczność z przyjętymi założeniami
i znanymi faktami.
\end{przyklad}

\section{Co wiemy przed wprowadzeniem zbiorów?}
\label{sec:logic-to-sets}
Potrafimy już czytać formuły z negacją, koniunkcją,
alternatywą, implikacją i równoważnością. Potrafimy
odróżnić $\forall$ od $\exists$, negować oba
kwantyfikatory, rozpoznawać zakresy zmiennych,
a także budować dowody bezpośrednie, przez
przeciwstawienie, nie wprost, przez przypadki
i przez indukcję. Te narzędzia są niezbędne,
aby precyzyjnie zdefiniować \emph{zawieranie zbiorów},
\emph{równość zbiorów}, \emph{relacje} i
\emph{funkcje}. Od tego rozpoczniemy następną część.
